\begin{activity}{From Conservation of Energy to Finite Differences}

\section*{Purpose}
To understand finite difference heat flow equations as direct consequences of energy conservation applied to small volumes of rock.

Consider a vertical column of sediment divided into small layers. Heat flows between layers, and each layer can gain or lose energy depending on its neighbours and any internal heat production.

Rather than solving a continuous equation, we will track how heat moves between
\emph{discrete points}.

\section*{Part A: Physical intuition}

\begin{itemize}
    \item If a layer is hotter than the ones above and below it, does it gain or lose heat?
    \item If heat is being produced inside the layer, what happens over time?
\end{itemize}

\section*{Part B: Setting up the system}

\begin{itemize}
    \item Sketch a one-dimensional grid of nodes labeled $T_{i-1}$, $T_i$, and $T_{i+1}$.
    \item Indicate the spacing $\Delta z$.
\end{itemize}

\section*{Part C: Heat flux}

Heat flows according to Fourier's law:
\begin{equation}
    q_z = -k \pdv{T}{z}
\end{equation}

\textbf{Tasks:}
\begin{itemize}
    \item Write an expression for heat flowing \emph{into} node $i$.
    \item Write an expression for heat flowing \emph{out of} node $i$.
\end{itemize}

\bigskip

\section*{Part D: Energy conservation}

Energy conservation for a small volume gives:
\[
    \text{Rate of change} = \text{in} - \text{out} + \text{produced}
\]

\textbf{Task:}
Combine your expressions to derive:
\begin{equation}
    \pdv{T_i}{t} =
    \kappa \frac{T_{i+1} - 2T_i + T_{i-1}}{\Delta z^2}
    + \frac{A_i}{\rho c_p}
\end{equation}

\section*{Part E: Interpretation}

Explain in words:

\begin{itemize}
    \item What does the term $T_{i+1} - 2T_i + T_{i-1}$ represent?
    \item What role does heat production $A_i$ play?
    \item What happens if $A_i = 0$?
\end{itemize}

\bigskip

\textbf{Key Insight} \\
Finite difference equations are not numerical tricks—they are simply energy conservation applied locally.

\end{activity}